% Generated draft: requires mathematical and source review.
\documentclass[11pt]{article}
\usepackage[margin=25mm]{geometry}
\usepackage{fontspec}
\setmainfont{Latin Modern Roman}
\usepackage{amsmath,amssymb,mathtools,unicode-math}
\setmathfont{Latin Modern Math}
\usepackage{enumitem,needspace,xurl,hyperref}
\setlength{\parindent}{0pt}
\setlength{\parskip}{5pt}
\setlength{\emergencystretch}{3em}
\newcommand{\N}{\mathbb{N}}
\newcommand{\Z}{\mathbb{Z}}
\newcommand{\Q}{\mathbb{Q}}
\newcommand{\R}{\mathbb{R}}
\newcommand{\C}{\mathbb{C}}
\newcommand{\F}{\mathbb{F}}
\newcommand{\A}{\mathbb{A}}
\newcommand{\PP}{\mathbb{P}}
\newcommand{\E}{\mathbb{E}}
\newcommand{\eps}{\varepsilon}
\newcommand{\dd}{\,\mathrm{d}}
\begin{document}
\section*{UnsolvedMath ID: 30005489}
\textbf{Source problem code:} catalogue entry
\subsection{Definitions and mathematical statement}
Let \(n\) be the index occurring in the supplied source notation, and write
\[
\chi(n,0)\colon M_{\mathrm{Dol}}(n,0)\longrightarrow A_n
\]
for its degree-zero Hitchin morphism, with Dolbeault moduli space \(M_{\mathrm{Dol}}(n,0)\) and entire Hitchin base \(A_n\). Write \(\operatorname{IC}_{M_{\mathrm{Dol}}(n,0)}\) for the intersection complex with the coefficient convention intended by the source, and set
\[
K_n=R\chi(n,0)_*\operatorname{IC}_{M_{\mathrm{Dol}}(n,0)}.
\]
Here \(R\chi(n,0)_*\) denotes derived direct image. For a constructible complex \(F\), its support means the closed set
\[
\operatorname{Supp}(F)=\overline{\bigcup_{j\in\Z}\{a\in A_n:\mathcal H^j(F)_a\ne0\}}.
\]
The full-support assertion is that every nonzero direct summand of \(K_n\) has support equal to \(A_n\); explicitly,
\[
K_n\simeq F\oplus G,\qquad F\ne0
\quad\Longrightarrow\quad
\operatorname{Supp}(F)=A_n.
\]
The summands are taken in the constructible derived category intended by the source. In a setting where the decomposition theorem expresses \(K_n\) as a direct sum of shifted simple perverse sheaves, this says that every such nonzero summand has support \(A_n\).

\textbf{Conjecture (Full support; catalogue entry).} The complex \(K_n\) has full support on the entire Hitchin base \(A_n\), including its complement of the reduced Hitchin locus. Thus no nonzero direct summand is supported on a proper subvariety of \(A_n\).

\textbf{Unresolved setup.} The supplied record does not define the underlying curve, its genus, the ground field, the structure group, the stability convention for the moduli space, the coefficient convention for the intersection complex, or the admissible range of \(n\). Consequently, the conjecture above preserves the source's notation and target but remains conditional on its intended geometric setup; it does not assert an extension to arbitrary choices of those data. Degree zero is explicit. No coprimality hypothesis is stated for this conjecture.

\textbf{Supplied dated qualification.} The curation dated 2026-08-22 labels the problem partially solved and reports full support on the reduced Hitchin base in degree zero and for important families, while describing the assertion over the entire base as open. This is a supplied literature assessment, not an independently verified mathematical or current-status claim.
\subsection{Short English statement}
For the source's degree-zero Hitchin map, the conjecture asks whether every nonzero direct summand of the derived direct image of the intersection complex is supported on the entire Hitchin base. The geometric setup and range of \(n\) are missing from the supplied record. Its dated curation reports partial results, which do not settle the stated entire-base assertion.
\subsection{Sources}
\begin{itemize}
\item Local UnsolvedMath registry: \texttt{data/unsolvedmath/problems.json}, permanent ID \(30005489\), title \emph{Full Support for Singular Hitchin Systems}, source problem code \texttt{catalogue entry}; supplied complete source record from the Oberwolfach Reports Open Problems collection.
\item Supplied report citation: \emph{Komplexe Analysis -- Differential and Algebraic methods in Kähler spaces} (2023), Oberwolfach Reports, pp.~965--1030, DOI 10.4171/owr/2023/18. \url{https://doi.org/10.4171/owr/2023/18}. The record gives proposed year 2023 and names no individual proposer.
\item Supplied literature source: \emph{Combinatorial decomposition theorem}. \url{https://arxiv.org/abs/2209.00621}.
\item Supplied literature source: \emph{Hitchin fibrations are Ngô fibrations}. \url{https://arxiv.org/abs/2502.04966}.
\item Supplied literature source: \emph{OWR 2023 report}. \url{https://ems.press/content/serial-article-files/47013}.
\item Literature assessment checked at \texttt{2026-08-22}: the supplied curation labels the problem partially solved, reports full support on the reduced Hitchin base in degree zero and for important families, and describes the entire-base conjecture as open. This is supplied dated curation, not independently verified; the record itself requests independent verification of time-sensitive status and cited claims.
\item The supplied registry labels the statement \texttt{corrected\_verified} and says it was checked against the official Oberwolfach report. This records the registry's verification claim only; no source was independently visited or checked for this draft.
\end{itemize}
\end{document}
